\documentclass[12pt,a4paper]{article}
\usepackage[utf8]{inputenc}
\usepackage[spanish]{babel}
\usepackage{amsmath,amssymb,amsfonts}
\usepackage{geometry}
\usepackage{hyperref}
\geometry{margin=2.5cm}

\title{\textbf{Constantes Matemáticas en Aritmética}}
\author{Compendio de Constantes Fundamentales}
\date{\today}

\begin{document}

\maketitle

\section*{Introducción}
La aritmética es la rama de las matemáticas que estudia los números y las operaciones elementales hechas con ellos. En esta disciplina destacan constantes fundamentales que sirven como base para la teoría de números, la representación del conteo y la definición de las propiedades aritméticas básicas.

\section{Constante de Pitágoras ($\sqrt{2}$)}

\subsection*{1. Significado, Símbolo y Explicación}
La constante de Pitágoras se representa mediante el símbolo $\sqrt{2}$. Es la longitud de la diagonal de un cuadrado cuyo lado mide la unidad ($1$). Históricamente, fue la primera constante que la escuela pitagórica descubrió como un \textbf{número irracional}, demostrando que no puede expresarse como el cociente exacto de dos números enteros $p/q$.

\subsection*{2. Valor Típico en Ingeniería y Ciencia}
En cálculos habituales de ingeniería, arquitectura y ciencia aplicada, se utilizan las siguientes aproximaciones:
\begin{equation*}
\sqrt{2} \approx 1.4142 \quad \text{o} \quad 1.41421356
\end{equation*}

\subsection*{3. Valor Exacto a 15 Decimales}
\begin{equation*}
\sqrt{2} \approx 1.414213562373095
\end{equation*}

\subsection*{4. Valor Exacto a 100 Decimales}
Al ser un número irracional, no posee un número finito de decimales ni un patrón periódico. Su valor exacto extendido a 100 dígitos decimales es:
\begin{align*}
\sqrt{2} \approx \; & 1.414213562373095048801688724209698078569671875376948073176679737990732... \\
& ...47846210703885038753432764157273501
\end{align*}
\textbf{Margen de error:} Como constante matemática pura definida axiomáticamente, no tiene incertidumbre experimental ni margen de error intrínseco.

\vspace{0.8cm}
\hrule
\vspace{0.8cm}

\section{Constante de Teodoro ($\sqrt{3}$)}

\subsection*{1. Significado, Símbolo y Explicación}
Representada como $\sqrt{3}$, es la longitud de la diagonal principal de un cubo unitario o la altura de un triángulo equilátero de lado 2. Su irrealidad e irracionalidad fueron demostradas históricamente por Teodoro de Cirene mediante la Espiral de Teodoro.

\subsection*{2. Valor Típico en Ingeniería y Ciencia}
Utilizada frecuentemente en ingeniería eléctrica (sistemas trifásicos) y geometría:
\begin{equation*}
\sqrt{3} \approx 1.732 \quad \text{o} \quad 1.7320508
\end{equation*}

\subsection*{3. Valor Exacto a 15 Decimales}
\begin{equation*}
\sqrt{3} \approx 1.732050807568877
\end{equation*}

\subsection*{4. Valor Exacto a 100 Decimales}
\begin{align*}
\sqrt{3} \approx \; & 1.732050807568877293527446341505872366942805253810380628055806979451933... \\
& ...01690880003708114618675724857567563
\end{align*}
\textbf{Margen de error:} Valor matemático exacto (sin margen de error experimental).

\end{document}
